% Capítulo 3 (Principais Trabalhos Relacionados) — Slide 3 — Posicionamento da dissertação
\begin{frame}{Posicionamento da dissertação}
  \centering
  \resizebox{\linewidth}{!}{%
  \begin{tabular}{@{}lccccc@{}}
    \toprule
    \textbf{Trabalho} & \textbf{Base usada} & \textbf{CV/regime} & \textbf{Políticas} & \textbf{Multinível} & \textbf{Indicadores} \\
    \midrule
    Oroojlooyjadid et al.\ (2022) & Sintética & $<1$ & 3 & Sim & 2 \\
    Yuna et al.\ (2023) & Testes intermitentes & $>1{,}5$ & 2 & Não & 2 \\
    Liu et al.\ (2024) -- HAPPO & Sintética & $<1$ & 4 & Sim & 2 \\
    Zabraoui et al.\ (2025) -- GA-DQN & Real (Walmart M5) & $<1$ & 7 & Sim & 3 \\
    \midrule
    \rowcolor{accenteal!12}Esta dissertação (Exp.\ 1, PB) & \textbf{Real (PB)} & $2{,}5$--$5{,}4$ & \textbf{12} & Sim & \textbf{5} \\
    \rowcolor{accenteal!12}Esta dissertação (Exp.\ 2, BA) & \textbf{Real (BA)} & \textit{Lumpy} & \textbf{12} & Sim & \textbf{5 + testes} \\
    \bottomrule
  \end{tabular}}
  \vspace{0.6em}

  \small \fonte{Capítulo 3, Tabela 3.2 (subconjunto)}
  \note{
    Esta tabela posiciona quantitativamente a dissertação frente aos
    trabalhos mais próximos. A maior parte da literatura de RL para inventário
    ainda avalia poucas políticas em ambientes sintéticos de baixa variabilidade,
    CV menor que 1. A exceção relevante é Zabraoui et al. (2025), que usa a base
    real Walmart M5, mas ainda permanece em baixa variabilidade e com escopo menor
    de políticas e indicadores.
    Esta dissertação se diferencia em três eixos simultaneamente: usa dados
    reais de uma rede varejista brasileira, opera em regime de alta
    intermitência, e avalia doze políticas com cinco indicadores -- CTI, NS,
    TR, BE e FP -- além de validação estatística formal no segundo
    experimento. Essa combinação de escala, realismo de dados e
    abrangência de métricas é o que diferencia esta avaliação da literatura
    revisada.
  }
\end{frame}
